%美赛模板：正文部分

%模板中需要使用Xelatex编译两次，所以可以反复尝试编译
\documentclass[12pt]{article}  % 官方要求字号不小于 12 号，此处选择 12 号字体
% \linespread{1.1} %可选：设置行间距1.1倍
% 本模板不需要填写年份，以当前电脑时间自动生成
% 请在以下的方括号中填写队伍控制号
\usepackage[CONTROL_NUMBER_PLACEHOLDER]{easymcm2}  % 载入 EasyMCM2；auto-mm-writing 启动时从 run.yaml.team.control_number 替换
\problem{X}  % auto-mm-writing 会从 run.yaml.chosen_problem 替换
% \usepackage{mathptmx}  % 这是 Times 字体，中规中矩 
\usepackage{palatino}  % mathpazo 这palatino是 COMAP 官方杂志采用的更好看的 Palatino 字体，可替代以上的 mathptmx 宏包
\usepackage{pdfpages}
\usepackage{longtable} %长表格包
\usepackage{tabu}
\usepackage{threeparttable}
\usepackage{listings} %代码排版包
\usepackage{paralist} %压缩列表环境包（紧凑列表）
\usepackage{makecell} %表格换行包
\usepackage{multirow} %表格排版包
%\usepackage{fourier} %使用花里胡哨符号的宏包
%------------------------------------------代码排版包(待研究)
%    \usepackage{minted}  %代码排版包
%    \usepackage{caption} %浮动体包
%    \usepackage{xcolor}
\usepackage{matlab-prettifier} %MATLAB代码排版包
\usepackage{tcolorbox} %美丽盒子包
%    \tcbuselibrary{skins}
%    \tcbuselibrary{minted}
\tcbuselibrary{listings,skins,breakable,xparse}
%    \usemintedstyle{lovelace} %特定代码配色环境
%\tcbuselibrary{breakable}


\usepackage{pgfornament} %花样装饰包(待研究)
\usepackage[ruled,linesnumbered]{algorithm2e} %伪代码排版包
\usepackage{tikz} %绘图包
\usetikzlibrary{arrows.meta, positioning, calc}
\usepackage{wrapfig} %图文环绕包
\usepackage{float} %精确位置控制包
\usepackage{lettrine} %首字下沉包
%-------------------------------------------交叉引用实现包
\usepackage{hyperref}
\usepackage{cleveref}

\graphicspath{{img/}}          % 此处{img/}为相对路径，注意加上“/”
\let\itemize\compactitem{}
\let\enditemize\endcompactitem{}
%-------------------------------------------附录代码排版盒子（起）
 \tcbset{%
 mylist/.style = {%
  colframe = gray,
  colback = white,
  coltitle = red!50!yellow!3!white,
  colbacktitle = white,
  listing only,
  attach boxed title to top center = {yshift = -\tcboxedtitleheight/2},
  enhanced,
  drop fuzzy shadow, % shadow around listings
  left = 6.5mm, % distance between left rule and line number 
  breakable, % enable listing box to break by page
  %enhanced jigsaw, % box not being closed when broken by page
  fonttitle=\small\bfseries\color{black}, % customise font of listing title
  before skip=20pt plus 2pt, % vertical space between listings and text
  after skip=20pt plus 2pt,
 },
 example/.style 2 args = {%
  mylist,
  title = {Listing \thetcbcounter: #1},
  label = {#2},
 },
 }
 \newtcblisting[auto counter, number within = section, list inside = mcode]{matlab}[3][]{%
  listing options = {%
   style = Matlab-editor,
   numbers = left,
   numberstyle = \footnotesize\color{darkgray}\fontfamily{pcr}\selectfont,
   basicstyle = \footnotesize\fontfamily{pcr}\selectfont,
   breaklines = true,  % 启用自动换行
   breakatwhitespace = true,  % 在空格处换行
   postbreak = \mbox{\textcolor{gray}{$\hookrightarrow$}\space}  % 换行指示符
  },
  overlay = {\begin{tcbclipinterior}\fill[blue!15!white] (frame.south west) rectangle ([xshift=5.3mm]frame.north west);\end{tcbclipinterior}},
  example = {#2}{#3}, #1,
 }
 
 % Python代码盒子定义 - 与MATLAB配色统一
 \newtcblisting[auto counter, number within = section, list inside = pcode]{python}[3][]{%
  listing options = {%
   language = Python,
   numbers = left,
   numberstyle = \footnotesize\color{darkgray}\fontfamily{pcr}\selectfont,
   basicstyle = \footnotesize\fontfamily{pcr}\selectfont,
   keywordstyle = \color{darkblue}\bfseries,
   commentstyle = \color{darkgreen},
   stringstyle = \color{red!60},
   showstringspaces = false,
   emphstyle = \color{orange!60},
   emph = {self, cls, True, False, None},
   breaklines = true,  % 启用自动换行
   breakatwhitespace = true,  % 在空格处换行
   postbreak = \mbox{\textcolor{gray}{$\hookrightarrow$}\space}  % 换行指示符
  },
  overlay = {\begin{tcbclipinterior}\fill[blue!15!white] (frame.south west) rectangle ([xshift=5.3mm]frame.north west);\end{tcbclipinterior}},
  example = {#2}{#3}, #1,
 }
%-------------------------------------------附录代码排版盒子（止）

% AI报告专用代码高亮样式（无边框，支持跨页）
\lstdefinestyle{ai_python_style}{
    language = Python,
    numbers = none,  % 去掉行号
    basicstyle = \footnotesize\fontfamily{pcr}\selectfont,
    keywordstyle = \color{darkblue}\bfseries,
    commentstyle = \color{darkgreen},
    stringstyle = \color{red!60},
    showstringspaces = false,
    emphstyle = \color{orange!60},
    emph = {self, cls, True, False, None},
    breaklines = true,
    breakatwhitespace = true,
    postbreak = \mbox{\textcolor{gray}{$\hookrightarrow$}\space},
    frame = none,  % 无边框
    backgroundcolor = \color{white},
    xleftmargin = 0pt,
    xrightmargin = 0pt,
    framexleftmargin = 0pt,
    framexrightmargin = 0pt,
    framextopmargin = 0pt,
    framexbottommargin = 0pt
}

\lstdefinestyle{ai_matlab_style}{
    style = Matlab-editor,
    numbers = none,  % 去掉行号
    basicstyle = \footnotesize\fontfamily{pcr}\selectfont,
    breaklines = true,
    breakatwhitespace = true,
    postbreak = \mbox{\textcolor{gray}{$\hookrightarrow$}\space},
    frame = none,  % 无边框
    backgroundcolor = \color{white},
    xleftmargin = 0pt,
    xrightmargin = 0pt,
    framexleftmargin = 0pt,
    framexrightmargin = 0pt,
    framextopmargin = 0pt,
    framexbottommargin = 0pt
}

\lstnewenvironment{ai_python}[1][]
    {\lstset{style=ai_python_style, #1}}
    {}

\lstnewenvironment{ai_matlab}[1][]
    {\lstset{style=ai_matlab_style, #1}}
    {}

\input{new_command.tex} %载入新命令
%\newcommand{\upcite}[1]{\textsuperscript{\textsuperscript{\cite{#1}}}} %
\title{Paper Title Placeholder \\ Replace with the actual title}  % auto-mm-writing 在 Step 3 用 run.yaml.chosen_title 替换

% 如需要修改题头（默认为 MCM/ICM），请使用以下命令（此处修改为 MCM）
%\renewcommand{\contest}{MCM}


%--------------------------------------------文档开始
 \begin{document}

% 摘要内容 — auto-mm-writing 从 abstract_draft.md (最终 pass) 注入
\begin{abstract}

    % ABSTRACT PLACEHOLDER — replaced by auto-mm-writing Step 5 (three-pass abstract).
    % If you are reading this in a built PDF, the writing stage did not run.
    % Each task's paragraph must contain hard numbers per Rule 9 of integrity-rules.md.

    Problem-type-and-overall-approach: one sentence stating what this paper does.

    For \textbf{Task 1}: one short paragraph with model, algorithm, and headline number.

    For \textbf{Task 2}: one short paragraph with model, algorithm, and headline number.

    For \textbf{Task 3}: one short paragraph with model, algorithm, and headline number.

    Management / policy insight: one paragraph translating findings into a recommendation.

    \vspace{5pt}
    \textbf{Keywords: Term1; Term2; Term3; Term4; Term5}

\end{abstract}
\maketitle  % 生成 Summary Sheet

\tableofcontents  % 生成目录
% part_1_pre.tex — placeholder scaffold
% auto-mm-writing fills each section from upstream stage artifacts:
%   stage0_triage/contest_brief.md       → Background
%   inputs/problems/<X>.pdf (paraphrase) → Problem Restatement
%   stage2_solving/validation.md         → Our Work
%   stage1_modeling/assumptions.md       → Assumptions
%   stage1_modeling/notation.md          → Notation

% Chapter 1: Introduction
\section{Introduction}
\subsection{Background}
% PLACEHOLDER — auto-mm-writing replaces with a 1-2 paragraph background
% derived from stage0_triage/contest_brief.md and the problem statement.
% Do NOT paste the whole problem PDF here.
Background text goes here.

\subsection{Problem Restatement}
% PLACEHOLDER — auto-mm-writing lists each sub-question with a one-sentence brief.
\begin{itemize}
    \setlength{\parsep}{0ex}
    \setlength{\topsep}{2ex}
    \setlength{\itemsep}{1ex}
    \item \textbf{Task 1.} One-sentence restatement of sub-question 1.
    \item \textbf{Task 2.} One-sentence restatement of sub-question 2.
    \item \textbf{Task 3.} One-sentence restatement of sub-question 3.
    \item \textbf{Task 4.} One-sentence restatement of sub-question 4 (if any).
\end{itemize}

\subsection{Our Work}
% PLACEHOLDER — one paragraph per sub-question, with method and headline result.
% Reference figures via \Cref{fig:...}; figures themselves live in img/.
Our-work summary goes here.

% Optional: insert the model flowchart figure once it exists in img/
% \begin{figure}[H]
%     \centering
%     \includegraphics[width=0.88\textwidth]{fig-model-flowchart.pdf}
%     \caption{End-to-end framework.}
%     \label{fig:model-flowchart}
% \end{figure}

%--------------------------------- Assumptions
\section{Assumptions and Explanations}
% PLACEHOLDER — auto-mm-writing reads stage1_modeling/assumptions.md and
% renders the numbered list here. Every assumption has Justification + Affects + If-wrong.
\begin{enumerate}[\bfseries \textit{Assumption} 1:]
    \item \textbf{Assumption 1 — one-line claim.}\\
    \textbf{\textit{Explanation:}} 1-3 sentence justification with citation or data evidence.
    \item \textbf{Assumption 2 — ...}\\
    \textbf{\textit{Explanation:}} ...
\end{enumerate}

%--------------------------------- Notation
\section{Notation}
% PLACEHOLDER — auto-mm-writing reads stage1_modeling/notation.md and renders
% the single consolidated symbol table. Sets / Parameters / Decision variables /
% Auxiliary variables / Functions.
\begin{table}[H]
    \centering
    \caption{Symbol table.}
    \label{tab:notation}
    \begin{tabular}{l l l}
        \hline
        Symbol & Description & Domain / Units \\
        \hline
        $\cdot$ & Placeholder & TBD \\
        \hline
    \end{tabular}
\end{table}

% part_2_model.tex — placeholder scaffold
% auto-mm-writing populates each section from the matching upstream artifact.

%--------------------------------- Data preprocessing
\section{Data Preprocessing}
% PLACEHOLDER — drawn from stage0_triage/data_recon.md and how anomalies
% were handled in stage2_solving/pipeline.py. Include a data-flow figure
% (fig-data-flow.pdf in img/) if available.
Data preprocessing description.

%--------------------------------- Model build (one section per sub-question)
\section{Task 1: <Title from problem restatement>}
\subsection{Model formulation}
% PLACEHOLDER — derived from stage1_modeling/model.md. Show notation,
% objective(s), constraint groups (each with a 1-paragraph English gloss),
% any non-trivial derivations.
Model description, formulas, and derivations.

\subsection{Algorithm}
% PLACEHOLDER — prose summary of stage2_solving/pipeline.py for this task.
% Each algorithmic component must answer "what problem feature does this address?"
% per integrity-rules.md Rule 5.
Algorithm description; cite stage2_solving/runs/<exp_id>/ for the implementation.

\subsection{Results}
% PLACEHOLDER — from stage2_solving/validation.md. Reference figures.
Results paragraph with hard numbers.

% (Repeat \section{Task N: ...} for each sub-question.)

%--------------------------------- Sensitivity
\section{Sensitivity Analysis}
% PLACEHOLDER — one subsection per `to-sweep` assumption from
% stage1_modeling/assumptions.md, populated from stage2_solving/sensitivity.md.
Sensitivity discussion.

%--------------------------------- Validation
\section{Model Validation}
% PLACEHOLDER — from stage2_solving/validation.md. Baseline comparison,
% small-instance exact Gap (if applicable), ablation findings.
Validation discussion.

% part_3_conclusion.tex — placeholder scaffold
% auto-mm-writing populates from stage1_modeling/model.md "Known limitations"
% and stage2_solving/validation.md.

\section{Summary of Findings}
% PLACEHOLDER — one paragraph restating the headline findings per sub-question.
Summary of what was found.

\subsection{Key Results}
\begin{itemize}
    \item Headline result 1 (with hard number).
    \item Headline result 2 (with hard number).
    \item Headline result 3 (with hard number).
\end{itemize}

\section{Strengths and Weaknesses}
% PLACEHOLDER — honest discussion. Strengths the model can defend; weaknesses
% the writing stage acknowledges (small-sample caveat, etc.).

\subsection{Strengths}
\begin{itemize}
    \item Strength 1.
    \item Strength 2.
\end{itemize}

\subsection{Weaknesses}
\begin{itemize}
    \item Weakness / limitation 1.
    \item Weakness / limitation 2.
\end{itemize}

\section{Future Work}
% Optional. Most contest papers can omit or keep this brief.
Brief paragraph on natural extensions.

%--------------------------------- References
% auto-mm-writing emits references.bib from stage1_modeling/literature.md.
% Bibliography style provided by easymcm2.sty.
\bibliographystyle{plain}
\bibliography{references}

% Memo to DWTS Producers (1-2 pages as required by problem statement)

\begin{letter}{Memo to DWTS Producers}
\label{sec:memo}

\begin{center}
\rule{0.9\textwidth}{1.5pt}\\[0.2cm]
{\Large\bfseries MEMORANDUM}\\[0.2cm]
\rule{0.9\textwidth}{1.5pt}
\end{center}

\vspace{0.2cm}

\noindent
\begin{tabular}{@{}ll}
\textbf{TO:} & Executive Producers, \textit{Dancing with the Stars} \\
\textbf{FROM:} & Team \#2603956, MCM Data Analysis Group \\
\textbf{DATE:} & February 2026 \\
\textbf{RE:} & Recommendations on Combining Judge Scores and Fan Votes \\
\end{tabular}

\vspace{0.3cm}

\subsection*{Executive Summary}
We analyzed 34 seasons of DWTS data to address: \textbf{how should judge scores and fan votes be combined to produce fair and engaging outcomes?} Our analysis reconstructs hidden fan vote shares, compares aggregation rules, and examines controversial cases (Bobby Bones in Season~27). We recommend a hybrid policy balancing audience engagement with technical credibility.

\subsection*{Key Findings}
\begin{enumerate}
\setlength{\itemsep}{0.3ex}
    \item \textbf{Fan votes can be reliably inferred.} Using elimination-consistency constraints, we reconstructed weekly fan vote shares that reproduce 91\% of historical eliminations.
    \item \textbf{The percentage rule better reflects audience preference.} When the two rules disagree, percentage more often retains contestants with higher fan support.
    \item \textbf{Judge-save reduces controversy but shifts power.} Introducing judge-selected elimination prevents outcomes like Bobby Bones' victory, but transfers significant influence to judges.
    \item \textbf{Judges and fans weight attributes differently.} Judges emphasize progress and partner skill; fans respond more to popularity and celebrity identity.
\end{enumerate}

\subsection*{Recommendations}
\begin{enumerate}
\setlength{\itemsep}{0.3ex}
    \item[\textbf{R1.}] \textbf{Adopt a hybrid aggregation policy.}
    Use \textit{percentage-based} aggregation in early weeks (Weeks~1--6) to maximize audience engagement. Switch to \textit{judge-selected bottom-three elimination} at mid-season (Week~7 onward) to ensure technical quality influences final placements.

    \item[\textbf{R2.}] \textbf{Monitor fan--judge divergence in real time.} Large divergences (top-3 fan support but bottom-3 judge ranking) are early warning signs of potential controversy.

    \item[\textbf{R3.}] \textbf{Use transparency strategically.} Highlight ``divergence weeks'' to enhance narrative tension while maintaining trust.
\end{enumerate}

\subsection*{Why This Approach Works}
\begin{itemize}
\setlength{\itemsep}{0.2ex}
    \item \textbf{Audience satisfaction:} Fans see their votes matter early; late outcomes reflect both popularity and skill.
    \item \textbf{Technical credibility:} Avoids outcomes where low-scoring contestants win (Season~27).
    \item \textbf{Production simplicity:} Fits existing scoring infrastructure---no new voting technology required.
    \item \textbf{Narrative value:} Divergence between fans and judges creates natural storylines that can be highlighted rather than hidden.
\end{itemize}

\subsection*{Implementation Considerations}
The hybrid policy can be implemented immediately within the existing production framework. Early-season percentage aggregation requires no changes to current voting infrastructure. The mid-season transition to judge-selected elimination mirrors the existing dance-off format, where judges already make subjective decisions. We recommend announcing the rule structure at the season premiere to set viewer expectations and enhance transparency. Real-time divergence monitoring can be achieved using the same data pipelines that currently track votes and scores.

\vspace{0.2cm}
\noindent\textit{Respectfully submitted,}\\
\textbf{Team \#2603956}

\end{letter}


% 论文正文到此结束（含Memo），共25页；AI Report不计入正文页数
\label{MainLastPage} % 标记正文最后一页（供页眉 "Page X of Y" 使用）
\clearpage

% AI Report 单独开页：页码从1开始，总页数为1
\setcounter{page}{1}
\rhead{\small Page \thepage\ of 1}

% part_4_Appendix.tex — AI Use Report (mandatory for MCM/ICM 2024+)
% Does NOT count toward the 25-page main paper limit.
% auto-mm-writing drafts this from progress.jsonl events; the team reviews
% and edits before final build.

\clearpage
\section*{AI Use Report}
\addcontentsline{toc}{section}{AI Use Report}
\label{sec:ai_report}

\subsection*{Tools used}

\begin{tabular}{l l l}
    \hline
    Tool & Version / Model & Used for \\
    \hline
    Claude & <Opus 4.7 / latest> & High-level planning, draft prose, code review \\
    auto-mm skill & <git sha or version> & End-to-end orchestration of this contest run \\
    % Add other tools used (GitHub Copilot, Wolfram Alpha, image-gen models, etc.)
    \hline
\end{tabular}

\subsection*{Use cases}

\textbf{Use case 1 — <e.g. assumption list drafting>}\\
Tool: <which one>. Date: <when>.\\
Prompt summary: ``<short summary>''.\\
Output: <what the tool produced; how the team integrated it>.

\textbf{Use case 2 — <e.g. code skeleton>}\\
Tool: <which one>. Date: <when>.\\
Prompt summary: ``<short summary>''.\\
Output: <what the tool produced; how the team integrated it>.

\textbf{Use case 3 — <e.g. proofreading>}\\
Tool: <which one>. Used for: <scope>.

% Add more use cases as appropriate. Be honest. Reviewers do not penalize honesty.

\subsection*{Acknowledgment}

The team remains fully responsible for the correctness, originality, and
integrity of this submission. AI-generated content was reviewed, edited,
and verified by team members before inclusion.




\end{document}
